\chapter{Heart}

\section*{Patient focus}
This chapter explains the examination or clinical problem in accessible language. Patients should use it to prepare questions and understand the purpose, experience, benefits, and limitations of imaging; local instructions and the official report take priority.


\label{ch:ch42_heart}

Doctors can assess your heart without open surgery using several imaging techniques. \textbf{Echocardiography} (an ultrasound of the heart) is the most common first test---it is painless, uses no radiation, and can show how well the heart pumps, whether the valves are working, and whether there is fluid around the heart. \textbf{Cardiac CT} (also called CT coronary angiography or CTA) uses X-rays and contrast dye to look at the coronary arteries for blockages and to assess calcium buildup (calcium score). \textbf{Cardiac MRI} uses powerful magnets to provide the most detailed pictures of the heart muscle, including whether areas of the heart have been damaged by a previous heart attack.

Each test has a different strength: echocardiography is best for valves and real-time function, CTA is best for coronary arteries, and cardiac MRI is best for heart muscle characterisation and complex congenital disease. Your doctor will choose the right test for your condition.

---

\section{Imaging Options}
Echocardiography evaluates structure, valves, function, and hemodynamics. CT assesses coronary arteries, aorta, pericardium, and anatomy; MRI characterizes myocardium and function; nuclear studies assess perfusion or viability. Select according to the clinical question.

\begin{table}[h]
\centering
\caption{Cardiac imaging modalities}
\begin{tabular}{p{2.8cm} p{3.5cm} p{3.5cm} p{3.5cm}}
\hline
\textbf{Modality} & \textbf{Best for} & \textbf{Limitations} & \textbf{Key measurements} \\
\hline
Cardiac CT (CTA) & Coronary arteries, calcium score, anatomy & Radiation, contrast, arrhythmia artefact & Calcium score (Agatston), stenosis grade \\
Cardiac MRI & Myocardial characterisation, viability, function & Contraindications (devices, clips), time, cost & LVEF, LGE pattern, T1/T2 mapping \\
Nuclear perfusion (SPECT/PET) & Myocardial ischaemia, viability & Low spatial resolution, radiation & Reversible vs fixed defects, ejection fraction \\
Transoesophageal echo & Endocarditis, LA appendage thrombus, valve repair planning & Semi-invasive, sedation needed & Vegetation size, thrombus detection \\
Transthoracic echo & Valves, systolic function, effusion & Acoustic windows, operator-dependent & LVEF, valve gradients, RA/RV pressures \\
\hline
\end{tabular}
\end{table}

\begin{pearlbox}[Choosing the right cardiac test]
Ask yourself: Is the question about \textbf{coronary arteries}? $\rightarrow$ CTA. Is it about \textbf{valves or real-time function}? $\rightarrow$ Echo. Is it about \textbf{myocardial tissue} (viability, infiltration, inflammation)? $\rightarrow$ Cardiac MRI. Is it about \textbf{perfusion}? $\rightarrow$ Nuclear or stress echo.
\end{pearlbox}

---

\section{CT Coronary Angiography (CTA)}
CTA of the coronary arteries is performed with ECG-gating (prospective or retrospective) and IV contrast.

\subsection{Coronary Artery Calcium Score}
The calcium score (Agatston units) is derived from a non-contrast ECG-gated CT. It is a powerful predictor of cardiovascular risk:

\begin{table}[h]
\centering
\caption{Coronary artery calcium score interpretation}
\begin{tabular}{p{2.5cm} p{4cm} p{4cm}}
\hline
\textbf{Score} & \textbf{Burden} & \textbf{Clinical implication} \\
\hline
0 & None & Very low risk ($< 1\%$ 10-year event rate) \\
1--99 & Mild & Low risk \\
100--399 & Moderate & Intermediate risk; consider statin therapy \\
$\geq 400$ & Severe & High risk; high likelihood of significant stenosis \\
\hline
\end{tabular}
\end{table}

\subsection{CTA Stenosis Grading}
\begin{itemize}
  \item \textbf{Normal}: No plaque
  \item \textbf{Minimal stenosis}: $< 25\%$
  \item \textbf{Mild stenosis}: $25$--$49\%$
  \item \textbf{Moderate stenosis}: $50$--$69\%$
  \item \textbf{Severe stenosis}: $70$--$99\%$
  \item \textbf{Occlusion}: $100\%$
\end{itemize}

---

\section{Cardiac MRI (CMR)}

\subsection{Function and Volumes}
CMR is the reference standard for left ventricular ejection fraction (LVEF), volumes, and mass. Steady-state free precession (SSFP) cine sequences provide real-time assessment of wall motion, valves, and shunts.

\subsection{Late Gadolinium Enhancement (LGE)}
LGE is obtained 10--15 minutes after gadolinium administration and identifies myocardial scar or fibrosis. The pattern of LGE is diagnostic:

\begin{table}[h]
\centering
\caption{LGE patterns in cardiomyopathies}
\begin{tabular}{p{3.5cm} p{5cm} p{4.5cm}}
\hline
\textbf{Condition} & \textbf{LGE pattern} & \textbf{Significance} \\
\hline
Arrhythmogenic RV cardiomyopathy & Fibrofatty replacement of RV free wall & Diagnostic criterion (Task Force Criteria) \\
Cardiac sarcoidosis & Patchy mid-wall or epicardial, often basal septum & Diagnostic and guides immunosuppression \\
Dilated cardiomyopathy & Mid-wall septal fibrosis & Prognostic marker; increased arrhythmia risk \\
Hypertrophic cardiomyopathy & RV insertion points, mid-wall & Risk stratification for sudden cardiac death \\
Ischaemic cardiomyopathy & Subendocardial or transmural, in a coronary territory & Identifies infarction, guides viability assessment \\
Myocarditis & Epicardial or mid-wall, non-coronary distribution (inferolateral typical) & Lake Louise Criteria for diagnosis \\
\hline
\end{tabular}
\end{table}

\begin{pearlbox}[Ischaemic vs non-ischaemic LGE]
\textbf{Ischaemic} LGE always starts from the subendocardium and may extend transmurally, respecting coronary territories. \textbf{Non-ischaemic} LGE is typically mid-wall or epicardial and does NOT follow a coronary distribution. This distinction is critical for determining aetiology.
\end{pearlbox}

\subsection{Tissue Characterisation}
\begin{itemize}
  \item \textbf{T1 mapping}: Elevated native T1 indicates diffuse fibrosis, oedema, or infiltration. Extracellular volume (ECV) fraction quantifies diffuse fibrosis.
  \item \textbf{T2 mapping}: Elevated T2 indicates myocardial oedema (acute inflammation, ischaemia).
  \item \textbf{T2* mapping}: Reduced T2* ($< 20$ ms) indicates iron overload; critical in thalassaemia patients.
\end{itemize}

---

\section{Echocardiography Basics}
Key measurements from transthoracic echocardiography:
\begin{itemize}
  \item \textbf{LVEF} (Simpson's biplane): Normal $\geq 55\%$; reduced $< 50\%$ suggests cardiomyopathy or ischaemic dysfunction
  \item \textbf{Diastolic function}: E/A ratio, e' velocity, E/e' ratio ($> 14$ suggests elevated filling pressures)
  \item \textbf{Valve assessment}: Colour Doppler for regurgitation (graded mild/moderate/severe); continuous-wave Doppler for stenosis gradients
  \item \textbf{Pulmonary artery systolic pressure} (PASP): Estimated from tricuspid regurgitation velocity
\end{itemize}

---

\section{Pericardial Disease}
\begin{itemize}
  \item \textbf{Pericardial effusion}: Low-attenuation fluid around the heart on CT; anechoic on echo. If large or causing chamber compression (tamponade), urgent drainage is required.
  \item \textbf{Constrictive pericarditis}: Thickened, calcified pericardium ($> 3$ mm) with impaired diastolic filling. CT and CMR both demonstrate pericardial thickening; CMR shows diastolic septal bounce.
  \item \textbf{Pericardial cyst}: Well-defined, water-attenuation mass at the right cardiophrenic angle. Usually incidental and asymptomatic.
\end{itemize}

---

\section{Cardiomyopathies}
\begin{table}[h]
\centering
\caption{Cardiomyopathies: imaging features}
\begin{tabular}{p{3cm} p{3.5cm} p{3.5cm} p{3cm}}
\hline
\textbf{Type} & \textbf{Echo/CMR findings} & \textbf{LGE pattern} & \textbf{Key features} \\
\hline
ARVC & RV dilation, aneurysms, fibrofatty replacement & RV free wall & Task Force Criteria include imaging criteria \\
DCM & Dilated LV, spherical remodelling, reduced LVEF & Mid-wall septal & Most common cardiomyopathy overall \\
HCM & Asymmetric septal hypertrophy ($\geq 15$ mm), LVOT obstruction & RV insertion points & Most common inherited cardiomyopathy; sudden death risk \\
RCM & Normal/small cavities, biatrial dilation, restrictive filling & Variable & Amyloid: diffuse subendocardial LGE; elevated native T1 \\
\hline
\end{tabular}
\end{table}

---

\section{Interpretation}
Describe chamber size, systolic function, wall motion, valves, pericardium, coronary or great-vessel findings, and comparison studies. Quantitative values require correct technique, reference population, and clinical context.

\section{Safety}
Consider rhythm, renal function, contrast, pregnancy, device compatibility, and patient stability. Acute ischemia, tamponade, dissection, or severe ventricular dysfunction requires prompt clinical correlation and escalation.

\clearpage
